\documentclass[11pt,a4paper]{article}
\usepackage[T1]{fontenc}
\usepackage[utf8]{inputenc}
\usepackage[english]{babel}
\usepackage{amsmath,amssymb}
\usepackage[margin=2cm]{geometry}
\setlength{\parindent}{0pt}
\setlength{\emergencystretch}{3em}
\allowdisplaybreaks

\title{About \texttt{false\_relation\_statements.tex}}
\author{}
\date{}
\begin{document}
\maketitle
This note explains what the test data in \texttt{false\_relation\_statements.tex} is used for. That file is input for the engine in \texttt{main.c}; it is read by the test \texttt{test\_set\_theory.c} in this folder.

\begin{enumerate}
\item \textbf{Purpose.} The file holds 39 statements that are not true, about relations given by a rule, ``there is exactly one'' and membership in a set of numbers. It is an anti-regression test: the engine must never report any of them as proved. It protects against changes that make the engine prove more statements by reading the notation less strictly.

\item \textbf{What the statements are.} Each one is close to a true statement of \texttt{relations\_by\_rule.tex}. A correspondence claimed to be the graph of a map when some $x$ has no $y$, or two. A relation claimed to be symmetric, reflexive or transitive when it is not, or claimed not to be when it is. A pair of numbers outside the two sets claimed to be in the relation. Five statements are traps for the substitution that unfolds $(s, t) \in R$: with the rule $y = 2x$ the pair $(a + 1, 2a + 1)$ is in $R$ only if the brackets around $a + 1$ are forgotten; with the rule $\exists k\ (y = kx)$ the statement $\forall k\ ((k, 1) \in R)$ becomes true only if the $k$ of the statement is captured by the quantifier of the rule; with a relation whose first letter is $Z$ and the rule $\exists k \in \mathbb{Z}\ (Z + k = y)$, the last statement becomes true only if the $Z$ of $\mathbb{Z}$ is replaced together with the letter. The last one was found as a real error of the engine and is kept so that it cannot come back.

\item \textbf{How the statements are written.} Every property is written out as a formula, not named.

\item \textbf{Adding statements.} Any new statement must be false. A statement that the engine should refuse because it cannot be read in one way only, such as a quantifier without a set that nothing forces, is not false: those are in \texttt{tests/engine/test\_engine.c}, with the answer expected for each.

\item \textbf{How to run it.} From the top folder of the repository:
\begin{verbatim}
cc -std=c11 -O2 -Wall -Wextra -o /tmp/test_set_theory \
    tests/set_theory/test_set_theory.c && /tmp/test_set_theory
\end{verbatim}

\item \textbf{Expected result.} The line \texttt{PASS: 39 of 39 false statements rejected}. Today the engine also proves the negation of all of them, and prints the numbers that show it.
\end{enumerate}
\end{document}
